%! Solving Exponential Equations with a Common Base — Unit 7, Lesson 7.3 — Group Activity
\documentclass[10pt]{article}
\usepackage{algebra2-article}
\usepackage{algebra2-boxes}
\newcommand{\TallMath}[1]{$\displaystyle #1\rule[-1.4em]{0pt}{3.2em}$}
\newcommand{\lgrid}[4]{%
  \draw[step=1, linegray!40, very thin] (#1,#3) grid (#2,#4);
  \draw[->, semithick, charcoal] (#1,0)--(#2,0) node[below right, font=\scriptsize]{$x$};
  \draw[->, semithick, charcoal] (0,#3)--(0,#4) node[left, font=\scriptsize]{$y$};}
% #1=a  #2=ln(b)  #3=h  #4=k  #5=xmin  #6=xmax   (ln2 = 0.693147, ln3 = 1.098612)
\newcommand{\expgraph}[6]{\draw[very thick, forest, domain=#5:#6, samples=140, smooth]
  plot (\x,{(#1)*exp((#2)*((\x)-(#3)))+(#4)});}
\newcommand{\solline}[3]{\draw[goldacc, very thick] (#1,#3)--(#2,#3);}

\begin{document}

\pageheader{Unit 7, Lesson 7.3}{Group Activity}
\namepartnerperiod

\vspace{0.08in}

\begin{headlinebox}{forestbg}
\small\textbf{Make the Bases Match --- Then Say Why You're Allowed To.} Every equation below is solved
the same way: rewrite both sides over one base, then set the exponents equal. Your partner's job is to
stop you after Step 1 and ask the two questions that keep you honest: \textbf{are both sides really
powers of the same base?} and \textbf{did the exponent get multiplied, or did you accidentally add?}
\end{headlinebox}

\vspace{0.06in}

\begin{tcolorbox}[colback=white, colframe=black!40, title=\textbf{Tier R --- Remediate},
    fonttitle=\bfseries, arc=1.5mm, left=3mm, right=3mm, top=2mm, bottom=2mm, breakable]
\textbf{Part 1 --- Warm the toolbox.} Write each number as a power of the base named in the row.
\begin{center}
\small
\renewcommand{\arraystretch}{1.45}
\begin{tabular}{|l|c|c|c|c|c|}
\hline
\textbf{Base $2$} & $8$ & $16$ & $64$ & $\frac12$ & $\frac{1}{16}$ \\
\textbf{exponent} & \blank{1.0cm} & \blank{1.0cm} & \blank{1.0cm} & \blank{1.0cm} & \blank{1.0cm} \\
\hline
\textbf{Base $3$} & $9$ & $27$ & $\frac19$ & $1$ & $\frac{1}{81}$ \\
\textbf{exponent} & \blank{1.0cm} & \blank{1.0cm} & \blank{1.0cm} & \blank{1.0cm} & \blank{1.0cm} \\
\hline
\textbf{Base $5$} & $25$ & $125$ & $\frac{1}{25}$ & $1$ & $625$ \\
\textbf{exponent} & \blank{1.0cm} & \blank{1.0cm} & \blank{1.0cm} & \blank{1.0cm} & \blank{1.0cm} \\
\hline
\end{tabular}
\end{center}

\vspace{5pt}
\textbf{Part 2 --- Six equations.} Rewrite, match, solve. Use the space for the common-base step.
\begin{center}
\renewcommand{\arraystretch}{1.7}
\begin{tabularx}{\linewidth}{@{}l X c@{\hspace{2em}} l X c@{}}
\hline
\textbf{(a)} $2^{x}=64$ & \blank{2.6cm} & $x=$ \blank{1.2cm} &
\textbf{(b)} $5^{\,x-1}=25$ & \blank{2.6cm} & $x=$ \blank{1.2cm} \\
\textbf{(c)} $3^{x}=\frac19$ & \blank{2.6cm} & $x=$ \blank{1.2cm} &
\textbf{(d)} $4^{x}=64$ & \blank{2.6cm} & $x=$ \blank{1.2cm} \\
\textbf{(e)} $10^{\,2x}=1000$ & \blank{2.6cm} & $x=$ \blank{1.2cm} &
\textbf{(f)} $7^{\,x+2}=1$ & \blank{2.6cm} & $x=$ \blank{1.2cm} \\
\hline
\end{tabularx}
\end{center}

\vspace{3pt}
\textbf{Part 3 --- Prove two of them.} Substitute your answers back in and show the arithmetic.

\vspace{2pt}
\hspace*{1.2em} (b) \; $5^{\,\blacksquare-1}=$ \blank{4.5cm} \hfill (f) \; $7^{\,\blacksquare+2}=$
\blank{4.5cm}

\vspace{4pt}
Item (f) surprises people. Why does an exponent of $0$ turn out to be the answer, when there is no
visible power on the right-hand side? \blank{7.0cm}
\end{tcolorbox}

\begin{tcolorbox}[colback=white, colframe=black!40, title=\textbf{Tier A --- Approaching Proficiency},
    fonttitle=\bfseries, arc=1.5mm, left=3mm, right=3mm, top=2mm, bottom=2mm, breakable]
\textbf{Part 1 --- Both sides need work.} In every one of these, at least one side has to be rewritten
before the bases match. Show that step.
\begin{center}
\renewcommand{\arraystretch}{1.8}
\begin{tabularx}{\linewidth}{@{}l X c@{\hspace{2em}} l X c@{}}
\hline
\textbf{(a)} $9^{x}=27$ & \blank{2.4cm} & $x=$ \blank{1.2cm} &
\textbf{(b)} $8^{\,x+1}=16$ & \blank{2.4cm} & $x=$ \blank{1.2cm} \\
\textbf{(c)} $\left(\frac13\right)^{x}=81$ & \blank{2.4cm} & $x=$ \blank{1.2cm} &
\textbf{(d)} $25^{x}=125^{\,x-1}$ & \blank{2.4cm} & $x=$ \blank{1.2cm} \\
\textbf{(e)} $4^{\,3x}=8^{\,x+2}$ & \blank{2.4cm} & $x=$ \blank{1.2cm} &
\textbf{(f)} $6^{\,2x-1}=\frac{1}{36}$ & \blank{2.4cm} & $x=$ \blank{1.2cm} \\
\hline
\end{tabularx}
\end{center}

\vspace{4pt}
\textbf{Part 2 --- The same answers, read off a picture.} The curve is $y=3^{x}$. Each gold line is one
right-hand side.
\begin{center}
\begin{minipage}[c]{0.40\linewidth}
\centering
\begin{tikzpicture}[scale=0.40]
  \lgrid{-3}{3}{-3}{10}
  \begin{scope}
    \clip (-3,-3) rectangle (3,10);
    \expgraph{1}{1.098612}{0}{0}{-3}{3}
  \end{scope}
  \solline{-3}{3}{9}
  \solline{-3}{3}{1}
  \draw[redacc, dashed, very thick] (-3,-2)--(3,-2);
  \node[goldacc, font=\scriptsize, left] at (-1.4,9.6) {$y=9$};
  \node[goldacc, font=\scriptsize, left] at (-1.4,1.6) {$y=1$};
  \node[redacc, font=\scriptsize, left] at (-1.4,-2.6) {$y=-2$};
\end{tikzpicture}
\end{minipage}
\begin{minipage}[c]{0.56\linewidth}
\small
Read each solution off the graph, then confirm it algebraically in the middle column.
\renewcommand{\arraystretch}{1.6}
\begin{tabularx}{\linewidth}{@{}l X c@{}}
\hline
\textbf{Equation} & \textbf{Common base} & \textbf{$x=$} \\
\hline
$3^{x}=9$ & \blank{2.0cm} & \blank{1.2cm} \\
$3^{x}=1$ & \blank{2.0cm} & \blank{1.2cm} \\
$3^{x}=-2$ & \blank{2.0cm} & \blank{1.2cm} \\
\hline
\end{tabularx}
\vspace{3pt}
How many times can a horizontal line cross an exponential curve? \blank{2.2cm}
\end{minipage}
\end{center}

\vspace{2pt}
For the row with no answer, say what rules it out --- and name the feature of the graph you used.
\blank{7.5cm}

\vspace{5pt}
\textbf{Part 3 --- Error analysis.} Name the error, then give the correct solution.

\vspace{2pt}
\textbf{(a)} Devon solved $8^{x}=4$ like this: ``$\left(2^{3}\right)^{x}=2^{2}$, so $2^{\,3+x}=2^{2}$,
so $3+x=2$ and $x=-1$.'' \writelines{2}

\vspace{2pt}
\textbf{(b)} Priya solved $9^{x}=27$ like this: ``$3^{2}=9$ and $3^{3}=27$, so $x=3$.''
\writelines{2}

\vspace{2pt}
One of those two students would be helped most by rereading Warm-Up item 2. Which one, and why?
\blank{7.0cm}
\end{tcolorbox}

\begin{tcolorbox}[colback=white, colframe=black!40, title=\textbf{Tier E --- Extension},
    fonttitle=\bfseries, arc=1.5mm, left=3mm, right=3mm, top=2mm, bottom=2mm, breakable]
\textbf{Part 1 --- Radicals are exponents too} (Lesson 6.1). Convert first, then match.
\begin{center}
\renewcommand{\arraystretch}{1.8}
\begin{tabularx}{\linewidth}{@{}l X c@{\hspace{2em}} l X c@{}}
\hline
\textbf{(a)} $2^{x}=\sqrt[3]{2}$ & \blank{2.4cm} & $x=$ \blank{1.2cm} &
\textbf{(b)} $4^{x}=\sqrt{8}$ & \blank{2.4cm} & $x=$ \blank{1.2cm} \\
\textbf{(c)} $\left(\frac15\right)^{\,x-1}=\sqrt[3]{25}$ & \blank{2.4cm} & $x=$ \blank{1.2cm} &
\textbf{(d)} $2^{x}\cdot 8^{\,x-1}=32$ & \blank{2.4cm} & $x=$ \blank{1.2cm} \\
\hline
\end{tabularx}
\end{center}
Item (d) needs a rule the others did not. Which one, and where did you use it? \blank{6.5cm}

\vspace{6pt}
\textbf{Part 2 --- The fine print, tested.} The one-to-one property says: if $b^{m}=b^{n}$ with $b>0$
and $b\neq1$, then $m=n$. Both conditions are doing real work.
\begin{enumerate}[label=(\alph*), leftmargin=1.9em, itemsep=6pt]
  \item Find two \emph{different} numbers $m$ and $n$ with $1^{m}=1^{n}$. \; $m=$ \blank{1.4cm}
        \quad $n=$ \blank{1.4cm}
  \item So how many solutions does $1^{x}=1$ have? \blank{4.0cm}
  \item In Lesson 7.1 you were told $b\neq1$ because the graph would be a flat line. State the
        \emph{solving} reason for the same restriction, in one sentence. \writelines{2}
  \item Now the other condition. Explain why a base like $-2$ is not allowed, using
        $(-2)^{1/2}$ as your example. \writelines{2}
\end{enumerate}

\vspace{5pt}
\textbf{Part 3 --- A colony, and the first thing this method cannot do.} A bacterial colony starts at
$50$ cells and doubles every hour, so after $t$ hours there are
\[
  P(t) = 50\cdot 2^{\,t} \quad\text{cells.}
\]
\begin{center}
\renewcommand{\arraystretch}{1.3}
\begin{tabular}{|c|c|c|c|c|c|c|c|}
\hline
$t$ (hours) & $0$ & $1$ & $2$ & $3$ & $4$ & $5$ & $6$ \\ \hline
$P(t)$ & $50$ & $100$ & \blank{1.2cm} & \blank{1.2cm} & \blank{1.2cm} & \blank{1.2cm} & \blank{1.2cm} \\ \hline
\end{tabular}
\end{center}
\begin{enumerate}[label=(\alph*), leftmargin=1.9em, itemsep=6pt]
  \item When does the colony reach $3200$ cells? Set up $50\cdot2^{t}=3200$, divide first, then match
        bases. \blank{6.0cm}
  \item When does it reach $12{,}800$ cells? \blank{4.0cm}
  \item Now try $1000$ cells. Divide first: $2^{t}=$ \blank{1.6cm}. What goes wrong at the
        common-base step? \writelines{2}
  \item Use the table to trap the answer: the colony passes $1000$ cells between $t=$ \blank{1.0cm}
        and $t=$ \blank{1.0cm} hours.
  \item Exactly one of these two has \emph{no} solution; the other has a solution you cannot yet find.
        Say which is which, and give a different reason for each. \\[3pt]
        \hspace*{1.2em} $50\cdot2^{t}=1000$ \qquad\qquad $50\cdot2^{t}=0$ \writelines{3}
\end{enumerate}
\end{tcolorbox}

\end{document}
